\section{Incentives, learning and remaining questions}\label{sec:06-outlook-1}
Potential structure guarantees neither efficient selection nor suitable
information. The preserved feasibility example is fully symmetric and has
an exact potential whose global maximum is the worst-welfare pure equilibrium.
A separate three-player common-interest coordination game, paying
$\epsilon>0$ at all-zero, $1$ at all-one and zero otherwise, has precisely
two strict pure equilibria for $0<\epsilon<1$. Its pure-equilibrium price of
anarchy is $1/\epsilon$ and price of stability is $1$. This elementary control
shows no efficiency bound follows from full ordinary symmetry alone.
The claim does not extend unchanged to four players: balanced profiles add
weak zero-welfare equilibria.

Strict unilateral improvements terminate in finite potential games
\citep{monderer1996}; simultaneous best responses can cycle. Under uniform
asynchronous log-linear updates with weights $\lambda^{u_i}$, the finite
exact-potential model has stationary law proportional to $\lambda^\Phi$ by
detailed balance. Its high-$\lambda$ limit selects potential maximizers, not
necessarily welfare maximizers. Full assumptions, proof and an exact rational
chain check are kept in research; no general learning or LLM convergence
claim follows. The companion binary-potential candidate supplies another
possible source of potential existence, under its unchanged hypotheses.

These initial results support a coherent research programme, with no originality
clearance yet. The most promising next question adds decentralized private
communication, verification or incentive constraints to the exact Bayesian
architecture gap. A second lead optimizes genuine checkpoint tasks under
partial feedback and shared actuator costs. Information budgets must distinguish
scalar measurements, precision, samples, bits and computation. Neither a
standard averaging proof nor a small control example alone warrants a broad
new theory of agentic AI. Empirical language-model experiments remain separate
from the idealized agents analyzed here.
